\section{Metodologia}

\begin{frame}{EADRC aplicado ao punho}
  \begin{columns}[T]
    \begin{column}{0.48\textwidth}
      \begin{itemize}[<+->]
        \item Sistema linear: o ângulo do punho tem parcelas voluntária e involuntária,
          \begin{equation*}
            \theta_3 = \theta_{v_3} + \theta_{i_3}.
          \end{equation*}
        \item Torque de controle apenas no punho, $\mathbf{u} = [0\ \ 0\ \ u_3]^\top$:
          \begin{equation*}
            \ddot{\theta}_3 = b_0 u_3 + \zeta(\bm{\tau}, \bm{\theta}, \dot{\bm{\theta}}),
            \quad
            b_0 = \left[\mathbf{J}^{-1}\right]_{33}.
          \end{equation*}
        \item $\zeta$ inclui o acoplamento com ombro e cotovelo, tratado como parte da perturbação total ($n = 2$).
      \end{itemize}
    \end{column}
    \begin{column}{0.48\textwidth}
      \begin{itemize}[<+->]
        \item 
        Erro ideal, que vai a zero quando $u_3$ compensa $\tau_{i_3}$:
          \begin{equation*}
            e_3 = \bm{\theta_{v_3}} - \theta_3 = -\theta_{i_3}.
          \end{equation*}
        Na prática, $\theta_{v_3}$ é desconhecido e estimado,
          \begin{equation*}
            \hat{e}_3 = \bm{\hat{\theta}_{v_3}} - \theta_3,
          \end{equation*}
        \alert{e a qualidade da estimativa limita a supressão alcançável.}
      \end{itemize}
    \end{column}
  \end{columns}
\end{frame}

\begin{frame}{Malha de controle e estimadores de movimento voluntário}
  \centering
  % Closed loop of the error-based ADRC applied to the wrist joint
% (needs TikZ libraries positioning, calc and fit, and the CBA colors from cba2026.sty)
\begin{tikzpicture}[
    > = latex,
    font = \small,
    block/.style = {
        draw,
        thick,
        rounded corners = 2pt,
        align = center,
        minimum height = 0.95cm,
        inner sep = 4pt,
    },
    eadrc/.style = {block, draw = CBAPurple, fill = CBALightPurple, minimum width = 2.2cm},
    plant/.style = {block, draw = CBAOrange, fill = CBALightOrange},
    estimator/.style = {block, draw = CBAPetrol, fill = CBALightPetrol},
    junction/.style = {draw, thick, circle, minimum size = 0.45cm, inner sep = 0pt},
]
    % Blocks
    \node [junction] (sum) {};
    \node [eadrc, right = 1.2cm of sum] (eso) {ESO\\baseado em erro};
    \node [eadrc, right = 1.6cm of eso] (law) {Lei de\\controle};
    \node [plant, right = 1.6cm of law] (plant) {Membro superior\\(3 GdL)};
    \node [estimator] (estimator) at ($(eso)!0.5!(law) + (0, -1.55)$)
        {Estimador de movimento voluntário\\(ZPLP ou EBMFLC)};

    % EADRC frame
    \node [
        draw = CBAPurple,
        dashed,
        rounded corners = 4pt,
        inner xsep = 0.25cm,
        inner ysep = 0.2cm,
        fit = (eso) (law),
        label = {[CBAPurple, font = \small\bfseries]north:EADRC},
    ] (eadrc frame) {};

    % Auxiliary coordinates
    \coordinate (u branch) at ($(law.east)!0.45!(plant.west)$);
    \coordinate (y branch) at ($(plant.east) + (0.45, 0)$);
    \coordinate (output) at ($(plant.east) + (1.3, 0)$);

    % Forward path
    \draw [->, thick] (sum) -- (eso) node [midway, above] {$\hat{e}_3$};
    \draw [->, thick] (eso) -- (law) node [midway, above] {$\hat{\bm{z}}$};
    \draw [->, thick] (law) -- (plant) node [pos = 0.45, above] {$u_3$};
    \draw [->, thick] (plant) -- (output) node [at end, above left] {$\theta_3$};

    % Control signal fed back to the observer
    \fill (u branch) circle (1.5pt);
    \draw [->, thick] (u branch) |- ($(eadrc frame.south) + (0, -0.2)$) -| (eso.south);

    % Torques acting on the upper limb
    \draw [<-, thick] (plant.north) -- ++(0, 0.6) node [above] {$\bm{\tau}_i + \bm{\tau}_v$};

    % Measured wrist angle fed back through the estimator (+) and directly (-)
    \fill (y branch) circle (1.5pt);
    \draw [->, thick] (y branch) |- (estimator.east);
    \draw [->, thick] (estimator.west) -| (sum.south)
        node [pos = 0.25, above] {$\hat{\theta}_{v_3}$}
        node [at end, below right = -1pt] {\scriptsize $+$};
    \draw [->, thick] (y branch) |- ($(estimator.south) + (0, -0.45)$) -| ($(sum.west) + (-0.6, 0)$)
        -- (sum.west) node [at end, above left = -1pt] {\scriptsize $-$};
\end{tikzpicture}


  \pause
  \vspace{0.3cm}
  \begin{columns}[T]
    \begin{column}{0.48\textwidth}
      \begin{exampleblock}{ZPLP}
        \small
        Butterworth passa-baixas de 1ª ordem, corte em 5 Hz, filtrado nos dois sentidos (fase zero) sobre as amostras disponíveis até o instante atual.
      \end{exampleblock}
    \end{column}
    \begin{column}{0.48\textwidth}
      \pause
      \begin{exampleblock}{EBMFLC \cite{Qi:2024}}
        \small
        Combinação linear adaptativa de 100 senoides entre 0 e 12 Hz; a parcela de 0 a 4 Hz é tomada como movimento voluntário.
      \end{exampleblock}
    \end{column}
  \end{columns}
\end{frame}

\begin{frame}{Simulações de Monte Carlo}
  \begin{columns}[T]
    \small
    \begin{column}{0.58\textwidth}
      \begin{itemize}
        \item Runge-Kutta de 4ª ordem, $t \in [0, 6]$~s, passo de 1~ms.
        \item 100 execuções por controlador: a primeira com o modelo nominal \cite{Gebai:2018}, as demais com rigidezes reamostradas ($c_i \propto k_i$ acompanha).
        \item Mesma semente aleatória para todos os controladores: comparação pareada.
      \end{itemize}
      %
      \pause
      \vspace{0.2cm}
      \centering
      \footnotesize
\begin{tabular}{l c c c c}
    \toprule
    Rigidez [N$\cdot$m/rad] & Mín. & Nominal & Máx. & Variação \\
    \midrule
    $k_1$ (ombro)    & 150 & 180 & 210 & $\pm16\%$ \\
    $k_2$ (cotovelo) & 50  & 70  & 90  & $\pm28\%$ \\
    $k_3$ (bíceps)   & 30  & 40  & 50  & $\pm25\%$ \\
    $k_4$ (punho)    & 5   & 10  & 15  & $\pm50\%$ \\
    \bottomrule
\end{tabular}

    \end{column}

    \pause
    \begin{column}{0.40\textwidth}
      \begin{itemize}
        \item Cenários: \alert{repouso} ($\bm{\tau}_v = \mathbf{0}$) e \alert{movimento} ($\bm{\tau}_v$ senoidal).
      \end{itemize}
      \centering
      \includegraphics[width=\linewidth]{figures/plots/uncontrolled_amplitude_1.0/torque_profiles_uncontrolled_amplitude_1.0.pdf}

      {\scriptsize Torques voluntários (movimento) e de tremor.}
    \end{column}
  \end{columns}
\end{frame}

\begin{frame}{Controladores e métricas de desempenho}
  \centering
  \footnotesize
\begin{tabularx}{\textwidth}{l >{\raggedright\arraybackslash}X l}
    \toprule
    Controlador & Parâmetros & Sintonia (modelo nominal) \\
    \midrule
    EADRC+ZPLP &
        $\omega_c = 10$, $\omega_o = 100$; Butterworth de 1ª ordem, corte em 5 Hz &
        Experimental, $\omega_o = 10\,\omega_c$ \\
    EADRC+EBMFLC &
        $\omega_c = 20$, $\omega_o = 80$; $n = 100$, $\mu = 0{,}005$, $T_p = 0{,}59$~s, $\alpha = 0{,}02$ &
        Experimental, $\omega_o = 4\,\omega_c$ \\
    IMC-PID &
        $K_p = 1{,}235$, $K_i = 1234{,}660$, $K_d = 0{,}506$ &
        IMC \cite{Rivera:1986}, punho isolado, $\lambda = 0{,}4\,\tau$ \\
    DE-PID &
        $K_p = 1{,}721$, $K_i = 15{,}059$, $K_d = 3{,}239$ &
        Evolução diferencial, $\theta_{v_3}$ conhecido \\
    \bottomrule
\end{tabularx}


  \pause
  \vspace{0.4cm}
  \begin{columns}[T]
    \begin{column}{0.48\textwidth}
      \begin{block}{Métricas}
        \small
        \begin{itemize}
          \item Resposta do punho: TPSR \cite{Qi:2024}, IAE \cite{Skogestad:2003}, $R^2$ \cite{Glantz:1990} e entropia \cite{Lathi:2018}.
          \item Sinal de controle: variação total (TV) \cite{Skogestad:2003} e potência \cite{Lathi:2018}.
        \end{itemize}
      \end{block}
    \end{column}
    \begin{column}{0.48\textwidth}
      \pause
      \begin{block}{Análise estatística}
        \small
        \begin{itemize}
          \item EADRC+ZPLP contra cada controlador: teste de Wilcoxon pareado \cite{Conover:1980}, com correlação bisserial de postos $r$.
          \item Correção de Holm-Bonferroni nas seis métricas; significância em $0{,}05$.
        \end{itemize}
      \end{block}
    \end{column}
  \end{columns}
\end{frame}
